\documentclass[aspectratio=169]{beamer}

\usetheme{Madrid}
\usecolortheme{default}
\setbeamertemplate{navigation symbols}{}

\usepackage{amsmath,amssymb,bm}
\usepackage{mathtools}

% ---------- Notation ----------
\newcommand{\x}{\bm{x}}
\newcommand{\uu}{\bm{u}}
\newcommand{\pp}{p}
\newcommand{\Rey}{\mathrm{Re}}
\newcommand{\avg}[1]{\overline{#1}}
\newcommand{\up}{\uu'}
\newcommand{\ubar}{\avg{\uu}}
\newcommand{\grad}{\nabla}
\newcommand{\Div}{\nabla\cdot}
\newcommand{\Lap}{\Delta}

\title[RANS II: Momentum]{RANS Part 2: Reynolds-Averaged Momentum Equation (x-component)}
\author{(adapted from Steve Brunton)}
\date{}

\begin{document}

% ------------------------------------------------------------
\begin{frame}
  \titlepage
\end{frame}

% ------------------------------------------------------------
\begin{frame}{Where we are (from Part 1)}
We start from incompressible Navier--Stokes (nondimensional):
\begin{align*}
\frac{\partial \uu}{\partial t}  + (\uu\cdot\grad)\uu &= -\grad \pp + \frac{1}{\Rey}\Lap\uu,\\
\Div\uu &= 0.
\end{align*}

Reynolds decomposition (time average):
\[
\uu(\x,t)=\ubar(\x)+\up(\x,t),\qquad \avg{\up}=0.
\]

\textbf{Result from Part 1:}
\[
\Div\ubar=0,\qquad \Div\up=0.
\]
\end{frame}

\begin{frame}{Other identities from part 1}
Recall we are taking averages in time. We know \\
\begin{align*}
\avg{u+v}=\avg{u}+\avg{v},\\
\avg{\avg{u}}=\avg{u},\\
\avg{\avg{u}\cdot v}=\avg{u}\cdot \avg{v},\\
\avg{\frac{\partial u}{\partial x}}=\frac{\partial \avg{u}}{\partial x}
\end{align*}

In general, 
\begin{align*}
\avg{u\cdot v}\neq \avg{u}\cdot \avg{v},\\
\avg{u'^2}\neq 0.
\end{align*}
\end{frame}
% ------------------------------------------------------------
\begin{frame}{Motivation: why average the momentum equation?}
Engineers often want \textbf{mean} quantities:
\begin{itemize}
  \item mean drag / mean lift,
  \item mean separation behavior,
  \item bulk transport trends,
\end{itemize}
without resolving every turbulent eddy in space-time.

\begin{block}{Plan today}
Plug $\uu=\ubar+\up$ (and $\pp=\avg{\pp}+\pp'$) into momentum,
then time-average to get an equation for $\ubar$.
\end{block}
\end{frame}

% ------------------------------------------------------------
\begin{frame}{We will derive only the x-momentum equation}
Write $\uu=(u,v,w)^T$. The x-component of momentum is:
\[
\frac{\partial u}{\partial t} + u\,\frac{\partial u}{\partial x} + v\,\frac{\partial u}{\partial y} + w\,\frac{\partial u}{\partial z}
= -\frac{\partial p}{\partial x} + \frac{1}{\Rey}\Lap u.
\]

We substitute:
\[
u=\avg{u}+u',\quad v=\avg{v}+v',\quad w=\avg{w}+w',\quad p=\avg{p}+p'.
\]
\end{frame}

% ------------------------------------------------------------
\begin{frame}{Averaging identities we will use repeatedly}
Time average operator $\avg{\cdot}$ is linear and (under mild regularity)
commutes with spatial derivatives:
\[
\avg{a+b}=\avg{a}+\avg{b},\qquad
\avg{\frac{\partial}{\partial x} a}=\frac{\partial}{\partial x} \avg{a}.
\]

Key facts:
\[
\avg{u'}=0,\qquad \avg{p'}=0,\qquad \avg{\Lap u'}=\Lap\avg{u'}=0.
\]

\begin{block}{Critical warning (the source of Reynolds stresses)}
In general,
\[
\avg{ab}\neq \avg{a}\,\avg{b}.
\]
So averages of \emph{products of fluctuations} need not vanish.
\end{block}
\end{frame}

% ------------------------------------------------------------
\begin{frame}{Step 1: substitute decomposition into x-momentum}
Start with
\[
\frac{\partial u}{\partial t} + u\,\frac{\partial u}{\partial x} + v\,\frac{\partial u}{\partial y} + w\,\frac{\partial u}{\partial z}
= -\frac{\partial p}{\partial x} + \frac{1}{\Rey}\Lap u.
\]

Substitute $u=\avg{u}+u'$, $v=\avg{v}+v'$, $w=\avg{w}+w'$, $p=\avg{p}+p'$.
Then
\[
\frac{\partial} {\partial t}(\avg{u}+u') + (\avg{u}+u')\frac{\partial}{\partial x}(\avg{u}+u')
+ (\avg{v}+v')\frac{\partial}{\partial y}(\avg{u}+u')
+ (\avg{w}+w')\frac{\partial}{\partial z}(\avg{u}+u')
= -\frac{\partial}{\partial x}(\avg{p}+p') + \frac{1}{\Rey}\Lap(\avg{u}+u').
\]
\end{frame}

% ------------------------------------------------------------
\begin{frame}{Step 1a: expand the convection term (x-part shown)}
The first convection piece expands like multiplying polynomials:
\[
(\avg{u}+u')\frac{\partial}{\partial x}(\avg{u}+u')
=
\avg{u}\,\frac{\partial \avg{u}}{\partial x}
+\avg{u}\,\frac{\partial u'}{\partial x}
+u'\,\frac{\partial \avg{u}}{\partial x}
+u'\,\frac{\partial u'}{\partial x}.
\]

Similarly,
\[
(\avg{v}+v')\frac{\partial}{\partial y}(\avg{u}+u')
=
\avg{v}\,\frac{\partial \avg{u}}{\partial y}
+\avg{v}\,\frac{\partial u'}{\partial y}
+v'\,\frac{\partial \avg{u}}{\partial y}
+v'\,\frac{\partial u'}{\partial y},
\]
and likewise for the $w$-term.
\end{frame}

% ------------------------------------------------------------
\begin{frame}{Step 2: time-average the entire equation}
Time-average term-by-term.

\textbf{Time derivative:}
\[
\avg{\frac{\partial}{\partial t}(\avg{u}+u')}
=
\frac{\partial}{\partial t} \avg{u} + \frac{\partial}{\partial t}\avg{u'}
=
\frac{\partial}{\partial t} \avg{u}.
\]
For a statistically stationary mean, $\frac{\partial}{\partial t}\avg{u}=0$.

\textbf{Pressure and viscosity:}
\[
\avg{-\frac{\partial}{\partial x}(\avg{p}+p')}=-\frac{\partial}{\partial x}\avg{p},\qquad
\avg{\Lap(\avg{u}+u')}=\Lap\avg{u}.
\]

\textbf{Mixed mean--fluctuation products die:} e.g.
\[
\avg{\avg{u}\,\frac{\partial u'}{\partial x}}=\avg{u}\,\avg{\frac{\partial u'}{\partial x}}=0,\quad
\avg{u'\,\frac{\partial \avg{u}}{\partial x}}=\avg{u'}\,\frac{\partial \avg{u}}{\partial x}=0.
\]
What survives are \emph{mean--mean} terms and \emph{fluctuation--fluctuation} terms.
\end{frame}

% ------------------------------------------------------------
\begin{frame}{Step 2 result: the averaged x-momentum equation (unsimplified)}
After averaging, the surviving terms are:
\begin{align*}
\avg{u}\,\avg{u}_x
+\avg{v}\,\avg{u}_y
+\avg{w}\,\avg{u}_z
&\;+\;\avg{u' u'_x}
+\avg{v' u'_y}
+\avg{w' u'_z} \\
&= -\avg{p}_x + \frac{1}{\Rey}\Lap \avg{u}.
\end{align*}

Note that we are using the standard shorthand where $u_x = \frac{\partial u}{\partial x}$, $u_y = \frac{\partial u}{\partial y}$, and $u_z = \frac{\partial u}{\partial z}$, and similar for velocities and pressure $v,w,p$.

\begin{block}{Interpretation}
Everything is now in terms of mean fields \emph{except} the three correlation terms
$\avg{u' u'_x}$, $\avg{v' u'_y}$, $\avg{w' u'_z}$.
These encode how fluctuations force the mean.
\end{block}
\end{frame}

% ------------------------------------------------------------
\begin{frame}{The problematic survivors: why they don't vanish}
Even though $\avg{u'}=0$, products of fluctuations need not average to zero:
\[
\avg{(u')^2}\neq 0,\qquad \avg{u'v'}\neq 0,\qquad \avg{u'w'}\neq 0.
\]

These are (components of) the \textbf{Reynolds stresses}.
They will become \textbf{spatial derivatives of} $\avg{(u')^2}$, $\avg{u'v'}$, $\avg{u'w'}$.

This is where closure enters: we need a model for these correlations.
\end{frame}

% ------------------------------------------------------------
\begin{frame}{Step 3: rewrite $\avg{u' u'_x}$ using a chain rule identity}
Use the product rule:
\[
\partial_x\big((u')^2\big)=2u'u'_x
\quad\Rightarrow\quad
u'u'_x=\frac12\,\partial_x\big((u')^2\big).
\]

Therefore,
\[
\avg{u'u'_x}
=
\frac12\,\frac{\partial}{\partial x}\avg{(u')^2}.
\]

(You can do the same trick for the other two terms, but one extra step is needed.)
\end{frame}

% ------------------------------------------------------------
\begin{frame}{Step 3: product-rule setup for the y- and z-terms}
Apply the product rule to $u'v'$ and $u'w'$:
\[
\frac{\partial}{\partial y}(u'v') = v' u'_y + u' v'_y,
\qquad
\frac{\partial}{\partial z}(u'w') = w' u'_z + u' w'_z.
\]

Rearrange:
\[
v' u'_y = \frac{\partial}{\partial y}(u'v') - u' v'_y,
\qquad
w' u'_z = \frac{\partial}{\partial z}(u'w') - u' w'_z.
\]

Now time-average:
\[
\avg{v'u'_y} = \frac{\partial}{\partial y}\avg{u'v'} - \avg{u'v'_y},\quad
\avg{w'u'_z} = \frac{\partial}{\partial z}\avg{u'w'} - \avg{u'w'_z}.
\]
\end{frame}

% ------------------------------------------------------------
\begin{frame}{The cancellation: use $\Div \up = 0$}
Collect the leftover terms:
\[
\avg{u'u'_x}+\avg{u'v'_y}+\avg{u'w'_z}.
\]

Factor out $u'$:
\[
u'u'_x + u'v'_y + u'w'_z
=
u'\,(u'_x+v'_y+w'_z)
=
u'\,(\Div\up).
\]

But from Part 1, $\Div\up=0$, hence
\[
\avg{u'u'_x}+\avg{u'v'_y}+\avg{u'w'_z}=0.
\]

So those leftover pieces vanish, and we keep only divergence of correlations.
\end{frame}

% ------------------------------------------------------------
\begin{frame}{Step 3 result: the Reynolds-stress-gradient form}
Using the identities and cancellations, the fluctuation terms become:
\[
\avg{u' u'_x}+\avg{v' u'_y}+\avg{w' u'_z}
=
\frac{\partial}{\partial x}\avg{(u')^2}
+\frac{\partial}{\partial y}\avg{u'v'}
+\frac{\partial}{\partial z}\avg{u'w'}.
\]

(Depending on your bookkeeping, you may see a factor of $\tfrac12$ appear and then cancel
when all three product-rule rearrangements are combined; the final packaged form is the spatial
divergence of the Reynolds-stress components shown above.)
\end{frame}

% ------------------------------------------------------------
\begin{frame}{Final x-component RANS momentum equation}
The Reynolds-averaged x-momentum equation can be written as:
\begin{align*}
\avg{u}\,\avg{u}_x
+\avg{v}\,\avg{u}_y
+\avg{w}\,\avg{u}_z
&=
-\avg{p}_x
+\frac{1}{\Rey}\Lap\avg{u}
-\Big(
\frac{\partial}{\partial x}\avg{(u')^2}
+\frac{\partial}{\partial y}\avg{u'v'}
+\frac{\partial}{\partial z}\avg{u'w'}
\Big).
\end{align*}

\begin{block}{Key interpretation}
The mean flow is driven by pressure + viscosity \emph{and} the divergence of Reynolds stresses.
\end{block}
\end{frame}

% ------------------------------------------------------------
\begin{frame}{Vector/tensor form (what this looks like in compact notation)}
Define the Reynolds-stress tensor:
\[
\bm{R} := \avg{\up\,\up^{T}}
=
\begin{bmatrix}
\avg{u'u'} & \avg{u'v'} & \avg{u'w'}\\
\avg{v'u'} & \avg{v'v'} & \avg{v'w'}\\
\avg{w'u'} & \avg{w'v'} & \avg{w'w'}
\end{bmatrix}.
\]

Then the mean momentum equation is commonly written:
\[
(\ubar\cdot\grad)\ubar
=
-\grad \avg{p}
+\frac{1}{\Rey}\Lap \ubar
-\Div \bm{R}.
\]

Along with mean incompressibility: \;\; $\Div\ubar=0$.
\end{frame}

% ------------------------------------------------------------
\begin{frame}{Why these are called ``stresses''}
The term $-\Div\bm{R}$ acts like an additional forcing (or stress divergence)
in the mean equation, with the same ``shape'' as viscous stress terms.

\begin{itemize}
  \item $\bm{R}$ depends on fluctuation correlations $\avg{u'_i u'_j}$.
  \item These are not determined by the mean equation alone.
\end{itemize}

\begin{block}{Closure problem (core of RANS modeling)}
The averaged equations introduce new unknowns ($\bm{R}$) that require modeling:
\[
\bm{R} \approx \mathcal{M}(\ubar, \grad \ubar, \ldots).
\]
\end{block}
\end{frame}

% ------------------------------------------------------------
\begin{frame}{What comes next (closure + ML)}
Next topics in a RANS sequence:
\begin{itemize}
  \item classic closures (e.g., eddy viscosity / Boussinesq hypothesis),
  \item $k$--$\varepsilon$, $k$--$\omega$, Spalart--Allmaras, etc.
  \item data-driven closure ideas (ML approximations of $\bm{R}$ or corrections).
\end{itemize}

\begin{block}{Takeaway}
RANS reduces space-time complexity, but pays a price: \textbf{closure modeling}.
\end{block}
\end{frame}

% ------------------------------------------------------------
\begin{frame}{Summary}
\begin{itemize}
  \item Start with x-momentum component of incompressible Navier--Stokes.
  \item Substitute $\uu=\ubar+\up$ and time-average.
  \item Mixed mean--fluctuation terms vanish; products of fluctuations remain.
  \item Use product-rule identities + $\Div\up=0$ to package surviving terms as
  \[
  -\left(\frac{\partial}{\partial x}\avg{(u')^2}+\frac{\partial}{\partial y}\avg{u'v'}+\frac{\partial}{\partial z}\avg{u'w'}\right)
  = -(\Div \bm{R})_x.
  \]
  \item The resulting mean equation requires a closure model for $\bm{R}$.
\end{itemize}
\end{frame}

\end{document}
